\section{Experiments}
\label{sec:experiments}

\paragraph{Pre-results status.}
This section fixes the evaluation questions, protocols, and table schemas; it
contains no new-method measurement or empirical conclusion. Red entries are
controlled evidence fields. Historical UCFRep dev84 diagnostics, proxy
readouts, and runs involving an untrained random Period Head are ineligible for
every table below.

\subsection{Evaluation Contract}
\label{sec:evaluation_contract}
\ClaimMap{C07,C10,C11,C12,C15}

\paragraph{Data and integrity.}
The primary benchmark will be the exact MultiRep release
\Protocol{multirep.release} under its official \Protocol{multirep.split}
protocol. The original MultiCounter paper specifies 1,157 synthetic videos,
52,590 periodic events, and a 7:2:1 train/validation/test split, while the
current MultiCounter+ repository reports the same size and provides a download
without a versioned release identifier or checksums
~\cite{tang2024multicounter,tang2026multicounterplus}. The final manifest must
therefore bind the release identifier, download source and time,
video/annotation/split SHA-256 digests, evaluator revision, and preprocessing
fingerprint. A run is table-ineligible if any binding is missing or if
incompatible artifacts are mixed. Before model comparison, a frozen dataset
characterization report must give split-level counts of videos, ground-truth
people, events, active-person levels, within-track tempo dispersion, detected
accelerations/decelerations, pauses, observation gaps, and identities with
enough cycles for every declared window length. Stratum definitions and counts
are fixed before looking at method errors.

\paragraph{Oracle and predicted tracks.}
The oracle-track protocol supplies official person identities and spatial
support, isolating repetition counting from detection and association. The
predicted-track protocol uses one frozen detector--pose--tracker frontend for
all compatible track-conditioned counters. Native joint MRAC systems retain
their audited official pipelines and appear in a separate contextual block,
not as if their inputs were controlled. Against ground-truth identities and
spatial support, HOTA and its DetA, AssA, and LocA components
~\cite{luiten2021hota} diagnose predicted-track detection, association, and
localization quality. IDF1 and ID-switch counts~\cite{ristani2016idmetrics}
provide complementary identity-duration and temporal-association diagnostics.
When the same frozen tracks feed several counters, these tracking values must
be identical and cannot be credited as counting gains.

\paragraph{Evaluator-ready outputs.}
Each counter adapter must emit the scored event tuples in
Equation~\ref{eq:event_schema}, the raw per-track count vector, and the exact
spatial support expected by the pinned evaluator. Oracle tracks retain official
identity support. Predicted tracks remain separate through evaluation: the
manifest-bound adapter must specify one deterministic matching policy and how
unmatched, duplicated, fragmented, merged, and identity-switched tracks become
false positives or false negatives. No post-hoc fragment joining is allowed
unless it is a separately audited method component. Adapter source, tests, and
per-instance outputs are required before Period-mAP or matched person-wise
count metrics become table-eligible. Predicted-track outputs are called
track-wise until this matching is applied.

\paragraph{Metrics and uncertainty.}
Using the frozen official evaluator, we will report Period-mAP, AP50, AP75,
AvgMAE, and AvgOBO, together with a matched person-count error
diagnostic after the evaluator's frozen matching step. Period-mAP under common
predicted tracks is the confirmatory metric; the primary contrast is our method
minus track-conditioned PAMS on identical videos and frontend outputs. A claim
of improvement requires the lower endpoint of its paired 95\% interval to be
above zero. Native-pipeline rows are descriptive and cannot establish
non-inferiority or superiority under a common-input protocol.

Every reproducible method uses at least three pre-registered seeds, reported
as mean and between-seed sample standard deviation. The confirmatory interval
uses one pre-declared paired hierarchical bootstrap: each replicate resamples
completed seed pairs and then videos within each sampled seed, retaining every
identity of a sampled video and recomputing Period-mAP from event records rather
than averaging per-video AP. Its point estimate is the mean of paired
seed-level differences. The manifest fixes the resample count, random seed,
failed-run policy, and exact numbers of completed seeds, videos, people, and
events. Secondary ablations are exploratory, with family-wise Holm correction
if formal significance statements are made. Official
single-checkpoint numbers remain explicitly marked and are excluded from
variance and significance claims. Test annotations are not used for tuning,
checkpoint selection, routing calibration, or early stopping.

The numerical stabilizer \(\epsilon\) is fixed before evaluation in the units
of each denominator. A sensitivity check spans representable smaller and larger
values and reports changes in response amplitude and count metrics by valid
support; it must expose, rather than hide, any coverage dependence introduced
by Equations~\ref{eq:valid_mean} and~\ref{eq:normalized_ola}.

\paragraph{Baselines and supervision.}
The mandatory native baselines are MultiCounter and MultiCounter+
~\cite{tang2024multicounter,tang2026multicounterplus}. Under the common track
interface we adapt RepNet~\cite{dwibedi2020repnet},
TransRAC~\cite{hu2022transrac}, PoseRAC~\cite{yao2025poserac}, and
PAMS~\cite{gao2026pams}; DeTRC~\cite{ma2026detrc} and
D$^2$-STX~\cite{wang2026d2stx} are candidate baselines and enter only if
runnable public implementations are available and their output semantics
pass the experiment audit. Tables disclose input modality, training labels,
track source, run count, adapter code, crop/pose construction, training split,
tuning budget, pretrained checkpoint, and output conversion. Comparisons across supervision regimes are
accuracy--annotation trade-offs, not claims of identical training information.

\subsection{Primary Questions}
\label{sec:primary_questions}
\ClaimMap{C03,C04,C05,C06,C12,C15}

The evaluation is organized around four questions. First, how does a
no-count-label pose counter compare on person-wise oracle-track MRAC and on
matched predicted tracks? Second, does local routing help beyond a track-conditioned
PAMS base when tempo changes within one identity? Third, does response-domain
overlap-add reduce boundary errors relative to summing independently decoded
windows? Fourth, how do tracking errors, active-person count, and upstream
preprocessing affect robustness and efficiency? Each question has a named
table or figure and an associated claim that remains withheld until the
evidence audit passes. The mechanism study must also compare against running
PAMS independently per track, sliding-window PAMS without a learned router,
deterministic tempo bins, non-overlapping and center-crop decoding, uniform
response averaging, and phase-aligned fusion.

\begin{table*}[t]
  \centering
  \small
  \setlength{\tabcolsep}{3.1pt}
  \caption{Planned MultiRep main results. Reproducible rows will show the
  mean and sample standard deviation over pre-registered seeds; official
  single-checkpoint rows will be marked. Protocol blocks answer different
  questions and will not be pooled into one ranking.}
  \label{tab:main_multirep}
  \begin{tabular}{lllccccc}
    \toprule
    Method & Track protocol & Count supervision &
    Period-mAP$\uparrow$ & AP50$\uparrow$ & AP75$\uparrow$ &
    AvgMAE$\downarrow$ & AvgOBO$\uparrow$ \\
    \midrule
    \multicolumn{8}{l}{\textit{Controlled oracle tracks}}\\
    RepNet & Oracle & \Protocol{sup.repnet} & \R{main.oracle.repnet.pmap} &
      \R{main.oracle.repnet.ap50} & \R{main.oracle.repnet.ap75} &
      \R{main.oracle.repnet.mae} & \R{main.oracle.repnet.obo} \\
    TransRAC & Oracle & \Protocol{sup.transrac} &
      \R{main.oracle.transrac.pmap} & \R{main.oracle.transrac.ap50} &
      \R{main.oracle.transrac.ap75} & \R{main.oracle.transrac.mae} &
      \R{main.oracle.transrac.obo} \\
    PoseRAC & Oracle & \Protocol{sup.poserac} & \R{main.oracle.poserac.pmap} &
      \R{main.oracle.poserac.ap50} & \R{main.oracle.poserac.ap75} &
      \R{main.oracle.poserac.mae} & \R{main.oracle.poserac.obo} \\
    PAMS & Oracle & \Protocol{sup.pams} & \R{main.oracle.pams.pmap} &
      \R{main.oracle.pams.ap50} & \R{main.oracle.pams.ap75} &
      \R{main.oracle.pams.mae} & \R{main.oracle.pams.obo} \\
    \textbf{Ours} & Oracle & \Protocol{sup.ours} & \R{main.oracle.ours.pmap} &
      \R{main.oracle.ours.ap50} & \R{main.oracle.ours.ap75} &
      \R{main.oracle.ours.mae} & \R{main.oracle.ours.obo} \\
    \midrule
    \multicolumn{8}{l}{\textit{Controlled common predicted tracks}}\\
    RepNet & Common predicted & \Protocol{sup.repnet} &
      \R{main.pred.repnet.pmap} & \R{main.pred.repnet.ap50} &
      \R{main.pred.repnet.ap75} & \R{main.pred.repnet.mae} &
      \R{main.pred.repnet.obo} \\
    TransRAC & Common predicted & \Protocol{sup.transrac} &
      \R{main.pred.transrac.pmap} & \R{main.pred.transrac.ap50} &
      \R{main.pred.transrac.ap75} & \R{main.pred.transrac.mae} &
      \R{main.pred.transrac.obo} \\
    PoseRAC & Common predicted & \Protocol{sup.poserac} &
      \R{main.pred.poserac.pmap} & \R{main.pred.poserac.ap50} &
      \R{main.pred.poserac.ap75} & \R{main.pred.poserac.mae} &
      \R{main.pred.poserac.obo} \\
    PAMS & Common predicted & \Protocol{sup.pams} & \R{main.pred.pams.pmap} &
      \R{main.pred.pams.ap50} & \R{main.pred.pams.ap75} &
      \R{main.pred.pams.mae} & \R{main.pred.pams.obo} \\
    \textbf{Ours} & Common predicted & \Protocol{sup.ours} & \R{main.pred.ours.pmap} &
      \R{main.pred.ours.ap50} & \R{main.pred.ours.ap75} &
      \R{main.pred.ours.mae} & \R{main.pred.ours.obo} \\
    \midrule
    \multicolumn{8}{l}{\textit{Native full-pipeline context; descriptive only}}\\
    MultiCounter & Native & \Protocol{sup.multicounter} &
      \R{main.native.multicounter.pmap} & \R{main.native.multicounter.ap50} &
      \R{main.native.multicounter.ap75} & \R{main.native.multicounter.mae} &
      \R{main.native.multicounter.obo} \\
    MultiCounter+ & Native & \Protocol{sup.multicounterplus} &
      \R{main.native.multicounterplus.pmap} &
      \R{main.native.multicounterplus.ap50} &
      \R{main.native.multicounterplus.ap75} &
      \R{main.native.multicounterplus.mae} &
      \R{main.native.multicounterplus.obo} \\
    \bottomrule
  \end{tabular}
\end{table*}

The quantitative interpretation of Table~\ref{tab:main_multirep} remains
withheld until every displayed binding passes the empirical audits.

\begin{table*}[t]
  \centering
  \small
  \setlength{\tabcolsep}{3.0pt}
  \caption{Required protocol disclosure accompanying the main table. Native
  rows retain their official pipelines and are not pooled with common-input
  rows. Every field is manifest-bound.}
  \label{tab:protocol_disclosure}
  \begin{tabular}{lllllll}
    \toprule
    Method & Result block & Input & Count labels & Period labels & Track source & Runs \\
    \midrule
    RepNet & Common & \Protocol{disc.repnet.input} & \Protocol{disc.repnet.count} &
      \Protocol{disc.repnet.period} & \Protocol{disc.common.track} & \Protocol{disc.repnet.runs} \\
    TransRAC & Common & \Protocol{disc.transrac.input} & \Protocol{disc.transrac.count} &
      \Protocol{disc.transrac.period} & \Protocol{disc.common.track} & \Protocol{disc.transrac.runs} \\
    PoseRAC & Common & \Protocol{disc.poserac.input} & \Protocol{disc.poserac.count} &
      \Protocol{disc.poserac.period} & \Protocol{disc.common.track} & \Protocol{disc.poserac.runs} \\
    PAMS & Common & \Protocol{disc.pams.input} & \Protocol{disc.pams.count} &
      \Protocol{disc.pams.period} & \Protocol{disc.common.track} & \Protocol{disc.pams.runs} \\
    \textbf{Ours} & Common & \Protocol{disc.ours.input} & \Protocol{disc.ours.count} &
      \Protocol{disc.ours.period} & \Protocol{disc.common.track} & \Protocol{disc.ours.runs} \\
    MultiCounter & Native context & \Protocol{disc.multicounter.input} &
      \Protocol{disc.multicounter.count} & \Protocol{disc.multicounter.period} &
      \Protocol{disc.multicounter.track} & \Protocol{disc.multicounter.runs} \\
    MultiCounter+ & Native context & \Protocol{disc.multicounterplus.input} &
      \Protocol{disc.multicounterplus.count} & \Protocol{disc.multicounterplus.period} &
      \Protocol{disc.multicounterplus.track} & \Protocol{disc.multicounterplus.runs} \\
    \bottomrule
  \end{tabular}
\end{table*}

\begin{table}[t]
  \centering
  \small
  \setlength{\tabcolsep}{2.7pt}
  \caption{Planned construction-path diagnostic under oracle tracks. LT: local
  tempo; SR: soft routing; BS: boundary stitching; TC: track-continuity
  objective. Causal conclusions use removals from the full model and the
  factorial controls described below, not this order-dependent path alone.}
  \label{tab:additive_ablation}
  \begin{tabular}{lcccccc}
    \toprule
    Variant & LT & SR & BS & TC & Period-mAP$\uparrow$ & AvgMAE$\downarrow$ \\
    \midrule
    Track-PAMS & & & & & \R{abl.base.pmap} & \R{abl.base.mae} \\
    $+$ LT & \checkmarkcell & & & & \R{abl.lt.pmap} & \R{abl.lt.mae} \\
    $+$ SR & \checkmarkcell & \checkmarkcell & & &
      \R{abl.sr.pmap} & \R{abl.sr.mae} \\
    $+$ BS & \checkmarkcell & \checkmarkcell & \checkmarkcell & &
      \R{abl.bs.pmap} & \R{abl.bs.mae} \\
    $+$ TC & \checkmarkcell & \checkmarkcell & \checkmarkcell &
      \checkmarkcell & \R{abl.tc.pmap} & \R{abl.tc.mae} \\
    \bottomrule
  \end{tabular}
\end{table}

The ablation interpretation remains withheld until all rows share one frozen
protocol and have the required paired records. The full mechanism study also
removes each component one at a time, crosses local-tempo estimation, routing,
and stitching in a small factorial design, matches single- and multi-expert
capacity, tests experts with and without direct \(\hat p\), and reports route
occupancy, entropy, collapse, and specialization by tempo stratum.

The appendix pre-registers router variants, expert counts, local-versus-global
tempo, shared batched execution versus independent same-checkpoint replicas,
boundary and continuity ablations,
tempo/occlusion/identity-switch stress, cross-dataset transfer, and latency,
memory, FLOPs, and throughput as the number of active identities grows. All
plots and tables must be generated from frozen JSON, CSV, or Parquet records;
manually transcribed values are inadmissible.
